\documentclass[10pt,a4paper]{amsart}
\usepackage{amsmath,amssymb,amsthm,graphicx}
% Retained minimal preamble for the Stacks Project chapter more-algebra.tex (More on Algebra).
% The chapter file pinned-original-source.tex carries only \input{preamble} on its first line, so no
% preamble text lives inside the chapter itself.  The Stacks Project uses one single shared preamble for
% every chapter; the snapshot stages it identically as sources/stacks-chow/preamble.tex and
% sources/stacks-etale-cohomology/preamble.tex.  The block below is copied mechanically, line for line,
% from sources/stacks-chow/preamble.tex lines 14-18 (the Xy-pic setup) and 154-217 (the native theorem
% declarations together with the mathematical macros).  Nothing was added, removed or redefined; the
% wrapper class and the amsmath/amssymb/amsthm/graphicx base packages are supplied by the assembler.
% Chapter: more-algebra.tex -- More on Algebra.
% For commutative diagrams we use Xy-pic
\usepackage[all]{xy}

% We use 2cell for 2-commutative diagrams.
\xyoption{2cell}

% Theorem environments.
%
\theoremstyle{plain}
\newtheorem{theorem}[subsection]{Theorem}
\newtheorem{proposition}[subsection]{Proposition}
\newtheorem{lemma}[subsection]{Lemma}

\theoremstyle{definition}
\newtheorem{definition}[subsection]{Definition}
\newtheorem{example}[subsection]{Example}
\newtheorem{exercise}[subsection]{Exercise}
\newtheorem{situation}[subsection]{Situation}

\theoremstyle{remark}
\newtheorem{remark}[subsection]{Remark}
\newtheorem{remarks}[subsection]{Remarks}

\numberwithin{equation}{subsection}

% Macros
%
\def\lim{\mathop{\mathrm{lim}}\nolimits}
\def\colim{\mathop{\mathrm{colim}}\nolimits}
\def\Spec{\mathop{\mathrm{Spec}}}
\def\Hom{\mathop{\mathrm{Hom}}\nolimits}
\def\Ext{\mathop{\mathrm{Ext}}\nolimits}
\def\SheafHom{\mathop{\mathcal{H}\!\mathit{om}}\nolimits}
\def\SheafExt{\mathop{\mathcal{E}\!\mathit{xt}}\nolimits}
\def\Sch{\mathit{Sch}}
\def\Mor{\mathop{\mathrm{Mor}}\nolimits}
\def\Ob{\mathop{\mathrm{Ob}}\nolimits}
\def\Sh{\mathop{\mathit{Sh}}\nolimits}
\def\NL{\mathop{N\!L}\nolimits}
\def\CH{\mathop{\mathrm{CH}}\nolimits}
\def\proetale{{pro\text{-}\acute{e}tale}}
\def\etale{{\acute{e}tale}}
\def\QCoh{\mathit{QCoh}}
\def\Ker{\mathop{\mathrm{Ker}}}
\def\Im{\mathop{\mathrm{Im}}}
\def\Coker{\mathop{\mathrm{Coker}}}
\def\Coim{\mathop{\mathrm{Coim}}}

% Boxtimes
%
\DeclareMathSymbol{\boxtimes}{\mathbin}{AMSa}{"02}

%
% Macros for moduli stacks/spaces
%
\def\QCohstack{\mathcal{QC}\!\mathit{oh}}
\def\Cohstack{\mathcal{C}\!\mathit{oh}}
\def\Spacesstack{\mathcal{S}\!\mathit{paces}}
\def\Quotfunctor{\mathrm{Quot}}
\def\Hilbfunctor{\mathrm{Hilb}}
\def\Curvesstack{\mathcal{C}\!\mathit{urves}}
\def\Polarizedstack{\mathcal{P}\!\mathit{olarized}}
\def\Complexesstack{\mathcal{C}\!\mathit{omplexes}}
% \Pic is the operator that assigns to X its picard group, usage \Pic(X)
% \Picardstack_{X/B} denotes the Picard stack of X over B
% \Picardfunctor_{X/B} denotes the Picard functor of X over B
\def\Pic{\mathop{\mathrm{Pic}}\nolimits}
\def\Picardstack{\mathcal{P}\!\mathit{ic}}
\def\Picardfunctor{\mathrm{Pic}}
\def\Deformationcategory{\mathcal{D}\!\mathit{ef}}

% --- appended 2026-10-12 (Chief): the project own definitions of the reference/slogan/history
% display environments, copied mechanically from the shared Stacks preamble lines 8-12
% (\usepackage{verbatim} plus the three \newenvironment lines, which the chapter preamble keeps
% verbatim from the project preamble). Nothing added, removed or redefined: these environments
% carry no mathematical content and the project itself typesets them as comments. Without these
% three lines a counted statement whose own text uses them fails with Environment slogan undefined.
\usepackage{verbatim}
\newenvironment{reference}{\comment}{\endcomment}
\newenvironment{slogan}{\comment}{\endcomment}
\newenvironment{history}{\comment}{\endcomment}
\usepackage{cleveref}
\crefname{originallemma}{original lemma}{original lemmas}
\Crefname{originallemma}{Original lemma}{Original lemmas}
\makeatletter
\newlabel{lemma-functorial}{{0624}{source}{}{sourcepointer-1}{}}
\newlabel{lemma-functorial@cref}{{[originallemma][0][]0624}{[1][0][]source}{}{}{}}
\makeatother
\crefname{originallemma}{original lemma}{original lemmas}
\Crefname{originallemma}{Original lemma}{Original lemmas}
\makeatletter
\newlabel{lemma-koszul-derived-completion-complete}{{091Y}{source}{}{sourcepointer-2}{}}
\newlabel{lemma-koszul-derived-completion-complete@cref}{{[originallemma][0][]091Y}{[1][0][]source}{}{}{}}
\makeatother
\crefname{originallemma}{original lemma}{original lemmas}
\Crefname{originallemma}{Original lemma}{Original lemmas}
\makeatletter
\newlabel{lemma-tensor-Rlim-exact}{{091K}{source}{}{sourcepointer-3}{}}
\newlabel{lemma-tensor-Rlim-exact@cref}{{[originallemma][0][]091K}{[1][0][]source}{}{}{}}
\makeatother
\crefname{originalsection}{original section}{original sections}
\Crefname{originalsection}{Original section}{Original sections}
\makeatletter
\newlabel{homology-section-truncations}{{0118}{source}{}{sourcepointer-4}{}}
\newlabel{homology-section-truncations@cref}{{[originalsection][0][]0118}{[1][0][]source}{}{}{}}
\makeatother
\title{Derived completeness is equivalent to being the derived limit of the tensor products with the Koszul system}
\author{The Stacks Project Authors}
\date{stacks-project:more-algebra.tex:local-snapshot-20261012; GNU FDL 1.2 or later}
\begin{document}
\maketitle
\setcounter{subsection}{1}
\noindent\textit{Editorial note.} All delivered mathematical text is the author's own, copied from the pinned Stacks Project chapter file more-algebra.tex; no statement or proof step was rewritten, reordered or removed. The counted statement is the labeled lemma environment at source lines 26341-26350 (canonical tag 091Z) and the counted proof is the proof environment at source lines 26352-26385, which immediately follows it and is its own proof.\par\medskip
\noindent\textit{Editorial note.} The counted statement is formulated inside the situation of the section; that situation environment is delivered as the context span of this card, with its boundaries taken exactly at the opening and closing lines of the environment, so the objects named in the statement can be read and the reference inside it resolves in the compiled document.\par\medskip
\noindent\textit{Editorial note.} The three stubs of this card cover the functoriality lemma cited inside the situation, the statement giving one implication and the exactness statement; their kinds are the type words appearing in the labels themselves, following the lane convention.\par\medskip
\noindent\textit{Editorial note.} The preamble span runs to line 90, the current end of the retained preamble, which includes the appended environment definitions; the counted text does not use them, and the span is taken whole so that the file version is unambiguous.\par\medskip
\noindent\textit{Editorial note.} The source URL is the source-level Stacks Project URL, as required when more than one card is produced from one source; the per-claim tag is carried by claim\_key and by this note. The counted statement and proof cite no bibliography entry, so this card needs no bibliography.\par\medskip
\noindent\textit{Allowed standard prerequisites (editorial).}\begin{itemize}
\item The derived limit of the tensor products with the Koszul complexes is derived complete with respect to the ideal, which gives one implication. invoked at proof lines 26353-26355. Disposes of the implication from the second condition to the first.
\item The construction sending an object to the derived limit of its tensor products with a system is exact in the object variable, so a filtration reduces the statement to its graded pieces. invoked at proof lines 26364-26368. Reduces the converse to the vanishing of the derived limits of the graded pieces of the filtration.
\item The graded pieces of the stupid filtration of the Koszul complex are direct sums of copies of the quotient by the ideal up to shifts, whose derived limits of tensor products are computed by the torsion functors. used at proof lines 26369-26384. Identifies the derived limits in question with direct sums of torsion functors, which vanish by the assumption of derived completeness.
\end{itemize}\medskip
\ifdefined\hypertarget\hypertarget{sourcepointer-1}{}\fi
\noindent\textit{Source reference.} Original lemma 0624: sources/stacks-more-algebra/pinned-original-source.tex, source lines 6612--6612. The reference occurs inside the situation environment delivered as the context span, where it produces the inverse system of Koszul complexes; it is a pointer needed only to define the objects of the statement.\par\medskip
\ifdefined\hypertarget\hypertarget{sourcepointer-2}{}\fi
\noindent\textit{Source reference.} Original lemma 091Y: sources/stacks-more-algebra/pinned-original-source.tex, source lines 26329--26329. The reference is cited at proof lines 26353-26355 for one implication; the converse, which is the substance of the proof, is inside the counted span.\par\medskip
\ifdefined\hypertarget\hypertarget{sourcepointer-3}{}\fi
\noindent\textit{Source reference.} Original lemma 091K: sources/stacks-more-algebra/pinned-original-source.tex, source lines 23915--23915. The reference is invoked at proof lines 26364-26366 to reduce the statement to the graded pieces of the filtration; the explicit computation of those graded pieces is inside the counted span.\par\medskip
\ifdefined\hypertarget\hypertarget{sourcepointer-4}{}\fi
\noindent\textit{Source reference.} Original section 0118: sources/stacks-chow/tags.txt, source lines 1354--1354. The reference points to the framework of stupid truncations used to build the filtration, cited at proof lines 26357-26358; it is a pointer to the construction and the use made of the filtration is inside the counted span.\par\medskip
\begin{situation}
\label{situation-koszul}
Let $A$ be a ring. Let $I = (f_1, \ldots, f_r) \subset A$. Let
$K_n^\bullet = K_\bullet(A, f_1^n, \ldots, f_r^n)$
be the Koszul complex on $f_1^n, \ldots, f_r^n$
viewed as a cochain complex in degrees $-r, -r + 1, \ldots, 0$.
Using the functoriality of Lemma \ref{lemma-functorial}
we obtain an inverse system
$$
\ldots \to K_3^\bullet \to K_2^\bullet \to K_1^\bullet
$$
compatible with the inverse system
$H^0(K_n^\bullet) = A/(f_1^n, \ldots, f_r^n)$
and compatible with the maps $A \to K_n^\bullet$.
\end{situation}

\begin{lemma}
\label{lemma-characterize-derived-complete-Koszul}
In Situation \ref{situation-koszul}. Let $K \in D(A)$.
The following are equivalent
\begin{enumerate}
\item $K$ is derived complete with respect to $I$, and
\item the canonical map $K \to R\lim (K \otimes_A^\mathbf{L} K_n^\bullet)$
is an isomorphism of $D(A)$.
\end{enumerate}
\end{lemma}

\begin{proof}
If (2) holds, then $K$ is derived complete with respect to $I$
by Lemma \ref{lemma-koszul-derived-completion-complete}.
Conversely, assume that $K$ is derived complete with respect to $I$.
Consider the filtrations
$$
K_n^\bullet \supset
\sigma_{\geq -r + 1}K_n^\bullet \supset
\sigma_{\geq -r + 2}K_n^\bullet \supset \ldots \supset
\sigma_{\geq -1}K_n^\bullet \supset
\sigma_{\geq 0}K_n^\bullet = A
$$
by stupid truncations (Homology, Section \ref{homology-section-truncations}).
Because the construction $R\lim(K \otimes E)$ is exact in
the second variable (Lemma \ref{lemma-tensor-Rlim-exact})
we see that it suffices to show
$$
R\lim \left(
K \otimes_A^\mathbf{L}
(\sigma_{\geq p}K_n^\bullet/ \sigma_{\geq p + 1}K_n^\bullet)
\right) = 0
$$
for $p < 0$. The explicit description of the Koszul complexes above
shows that
$$
R\lim \left(
K \otimes_A^\mathbf{L}
(\sigma_{\geq p}K_n^\bullet/ \sigma_{\geq p + 1}K_n^\bullet)
\right) =
\bigoplus\nolimits_{i_1, \ldots, i_{-p}}
T(K, f_{i_1}\ldots f_{i_{-p}})
$$
which is zero for $p < 0$ by assumption on $K$.
\end{proof}

\end{document}
